\SourcePage{398}{401}
\section{Multiplet terms. Cases $c$ and $d$}

Besides coupling cases $a$ and $b$, and the cases intermediate between them, other coupling types also exist. Their origin is as follows. The appearance of the quantum

\SourcePage{399}{402}
number $\Lambda$ is ultimately due to the electric interaction of the two atoms in the molecule, which gives axial symmetry to the problem of determining the electronic terms (this interaction in a molecule is described as the coupling of the orbital angular momentum to the axis). The magnitude of this interaction is measured by the separations between terms having different values of $\Lambda$. Throughout the preceding discussion, this interaction was tacitly assumed to be so strong that these separations are large in comparison both with the multiplet-splitting intervals and with the intervals in the rotational structure of the terms. The reverse situation also occurs, however, when the interaction of the orbital angular momentum with the axis is comparable with, or even small in comparison with, the other effects. In such cases the projection of the orbital angular momentum on the axis cannot be regarded as conserved in any approximation, and the number $\Lambda$ therefore loses its meaning.

When the coupling of the orbital angular momentum to the axis is weak compared with the spin--orbit coupling, the system is said to be in \emph{case $c$}. This case occurs in molecules containing a rare-earth atom. Such atoms have $f$-electrons with uncompensated angular momenta; their interaction with the molecular axis is weakened because the $f$-electrons lie deep within the atom. Coupling types intermediate between $a$ and $c$ occur in molecules composed of heavy atoms.

When the coupling of the orbital angular momentum to the axis is weak compared with the intervals of the rotational structure, the system is said to be in \emph{case $d$}. This case is encountered for high rotational levels (with large $J$) of certain electronic terms of the lightest molecules ($\mathrm H_2$, $\mathrm{He}_2$). These terms contain a strongly excited electron whose interaction with the remaining electrons (or, as it is called, with the molecular ``core'') is so weak that its orbital angular momentum is not quantized along the molecular axis (whereas the core has a definite angular momentum $\Lambda_{\text{core}}$ about that axis).

As the distance $r$ between the nuclei increases, the interaction of the atoms becomes weaker and eventually becomes small in comparison with the spin--orbit interaction within the atoms. Consequently, electronic terms at sufficiently large $r$ belong to case $c$. This circumstance must be kept in mind when establishing the correspondence between molecular electronic terms and the atomic states obtained as $r\to\infty$.

In \S~80 we considered this correspondence while neglecting the spin--orbit interaction. When the fine structure of the terms is included, there is the additional question of the correspondence between the values $J_1$ and $J_2$ of the total angular momenta of the isolated atoms

\SourcePage{400}{403}
and the values of the molecular quantum number $\Omega$. We give the results here without repeating the reasoning, which is entirely analogous to that used in \S~80.

If the molecule consists of unlike atoms, the possible values of $\Omega$\footnote{When the two total atomic angular momenta $\mathbf J_1$ and $\mathbf J_2$ are combined into the resultant angular momentum $\boldsymbol\Omega$, the sign of $\Omega$ is clearly immaterial.} formed by joining atoms with angular momenta $J_1$ and $J_2$ ($J_1\geqslant J_2$) are determined by the same scheme (80.1), with $J_1,J_2$ in place of $L_1,L_2$ and $|\Omega|$ in place of $\Lambda$. The sole difference is that, when $J_1+J_2$ is half-integral, the smallest value of $|\Omega|$ is $1/2$ rather than the zero shown in the scheme. When $J_1+J_2$ is integral, there are $2J_2+1$ terms with $\Omega=0$, for which (as with $\Sigma$-terms when fine structure is neglected) their sign must be determined. If $J_1$ and $J_2$ are both half-integral, the number $2J_2+1$ is even, and there are equal numbers of terms conventionally denoted by $0^+$ and $0^-$. If $J_1$ and $J_2$ are both integral, there are $J_2+1$ terms $0^+$ and $J_2$ terms $0^-$ when $(-1)^{J_1+J_2}P_1P_2=1$, and the reverse when $(-1)^{J_1+J_2}P_1P_2=-1$.

If the molecule consists of identical atoms in different states, the resulting molecular states are the same as for unlike atoms, except that the total number of terms is doubled: each term occurs once as even and once as odd.

Finally, if the molecule consists of identical atoms in identical states (with angular momenta $J_1=J_2=J$), the total number of states remains the same as for unlike atoms, while their parity distribution is
\[
\begin{array}{lll}
\text{if }J\text{ is integral,}&\Omega\text{ even:}&N_g=N_u+1,\\
&\Omega\text{ odd:}&N_g=N_u,\\
\text{if }J\text{ is half-integral,}&\Omega\text{ even:}&N_u=N_g,\\
&\Omega\text{ odd:}&N_u=N_g+1.
\end{array}
\]
All the $0^+$ terms are then even and all the $0^-$ terms are odd.

As the nuclei approach one another, coupling of type $c$ usually passes into coupling of type $a$\footnote{The correspondence between the case-$a$ and case-$c$ classifications of terms cannot be established in general. It requires a specific examination of the potential-energy curves with allowance for the non-crossing rule for levels of the same symmetry (\S~79).}. The following noteworthy situation can then arise.

As stated above, a term with $\Lambda=0$ belongs to case $b$; from the standpoint of the case-$a$ classification this means that the multiplet levels having different values of $\Omega$ (and the same $\Lambda=0$) have the same energy. Such levels, however,

\SourcePage{401}{404}
may arise as atoms in different fine-structure states approach each other.

Thus, different pairs of atomic fine-structure states may correspond to the same molecular term. An analogous situation may occur for terms with $\Omega=0$ that, as the nuclei approach, pass into a molecular term with $\Lambda\ne0$ (and hence $\Sigma=-\Lambda$). Such levels are doubly degenerate because, in case $a$, the terms $0^+$ and $0^-$ (which may arise from different pairs of atomic states) have the same energy\footnote{We neglect here the so-called $\Lambda$-doubling (see \S~88).}.
